% !TeX program = xelatex

\documentclass[14pt, b4paper]{extreport}

\usepackage[
margin=1in
%top=1in,
%bottom=1in,
%left=0.5in,
%right=0.5in
]{geometry}

% For the title page
\usepackage{xcolor} 
\usepackage{graphicx}
\usepackage{ifthen}
\usepackage{tikz}
\usetikzlibrary{decorations.markings}
\usetikzlibrary{arrows.meta}
\usetikzlibrary{hobby, bending}
\tikzset{
	mydash/.style={dashed, dash pattern=on 6pt off 2pt},
	arrow inside/.style 2 args = {
		postaction={
			decorate,
			decoration={
				markings,
				mark=at position {#1*\pgfdecoratedpathlength-0.15cm} with {\coordinate (bent arrow 1);},
				mark=at position {#1*\pgfdecoratedpathlength-0.05cm} with {\coordinate (bent arrow 2);},
				mark=at position {#1*\pgfdecoratedpathlength+0.05cm} with {\coordinate (bent arrow 3);},
				mark=at position {#1*\pgfdecoratedpathlength+0.15cm} with {
					\coordinate (bent arrow 4);
					\ifthenelse{\equal{#2}{for}}
					{\draw[-{Triangle[length=3mm, width=1.6mm, bend]}] 
						(bent arrow 1) to[curve through={(bent arrow 2) .. (bent arrow 3)}] (bent arrow 4);
					}
					{
						\draw[-{Triangle[length=3mm, width=1.6mm, bend]}] 
						(bent arrow 4) to[curve through={(bent arrow 3) .. (bent arrow 2)}] (bent arrow 1);
					}
				}
			}
		}
	}
}

\def\ANS#1{\hspace*{2em}\textit{ج}: #1}
\def\SUG#1{\hspace*{2em}\textit{اقتراح}: #1}

%\usepackage{indentfirst}
\usepackage{float}

\newcounter{myfig}
\makeatletter
\newcommand{\myfig}[1][]{
	\refstepcounter{myfig}
	\textbf{الشكل \themyfig}
	\ifx\@empty#1\@empty
	\else
	\\[-2pt]#1
	\fi
}
\makeatother
\newcommand{\FIGURE}[1]{\noindent\hspace{20pt}\parbox[b]{3cm}{\myfig[#1]}\\}


% For changing the headers formatting
\usepackage{titletoc, titlesec}
\usepackage[inline]{enumitem}
\setlist{label=\textbf{\arabic*.}}
%\settasks{label-offset=25pt,label=(\xslalph*)}
\usepackage{environ}
\NewEnviron{tasks}
{
	\begin{enumerate*}[label=(\xslalph*), itemjoin=\hfil]
		\BODY
	\end{enumerate*}
}

\NewEnviron{tasks*}
{
	\begin{enumerate*}[
		before=\noindent,
		topsep=0pt,
		partopsep=0pt,
		parsep=5pt,
		itemsep=0pt,
		itemjoin=\\,
		label=(\xslalph*)\,\,,
		mode=unboxed
		]
		\BODY
	\end{enumerate*}
}

\makeatletter
\newcommand{\xslalph}[1]{\expandafter\@xslalph\csname c@#1\endcsname}
\newcommand{\@xslalph}[1]{%
	\ifcase#1
	\or أ%
	\or ب%
	\or جـ%
	\or د%
	\or هـ%
	\or و%
	\or ز%
	\or حـ%
	\or ط%
	\or ي%
	\or ك%
	\or ل%
	\or م%
	\or ن%
	\or س%
	\or ع%
	\or ف%
	\or ص%
	\or ق%
	\or ر%
	\or ش%
	\or ت%
	\or ث%
	\or خـ%
	\or ذ%
	\or ض%
	\or ظ%
	\or غ%
	\else \@ctrerr
	\fi
}
\AddEnumerateCounter{\xslalph}{\@xslalph}{m}
\makeatother

% For math
\usepackage{amsmath, amssymb, amsthm, needspace, amsfonts}

\newcommand{\conj}[1]{\mkern 0.5mu\overline{\mkern-0.5mu#1\mkern-0.5mu}\mkern 0.5mu}

\allowdisplaybreaks

% To use the arabic language
\usepackage{polyglossia}

\usepackage[]{setspace}
\setstretch{2}

\usepackage{tocbasic}
\addtotoclist[report.cls]{toc}
\renewcommand{\tableofcontents}{\listoftoc[{\contentsname}]{toc}}
\AfterTOCHead[toc]{\thispagestyle{empty}\pagestyle{empty}}
\AfterStartingTOC[toc]{\clearpage}


%\renewcommand{\thepage}{{\arabic{page}}}


% For fancy headers
\usepackage[twoside]{fancyhdr}
\pagestyle{fancy}
\fancyhf{}
\fancyhead[RO]{\ar{\leftmark\qquad\qquad\textbf{\thepage}}}
\fancyhead[RE]{\ar{\rightmark\qquad\qquad\textbf{\thepage}}}
\fancyhead[LO]{\small\ar{الفصل \thechapter}}
\fancyhead[LE]{\small{\ar{\textbf{بــنــد \arabic{section}}}}}

\renewcommand{\chaptermark}[1]{\markboth{#1}{}}
\renewcommand{\sectionmark}[1]{\markright{#1}{}}
\renewcommand{\headrulewidth}{0pt}
\setlength{\headsep}{3em}

\fancypagestyle{plain}{
	\fancyhf{}
	\fancyfoot[L]{\textbf{\thepage}}
}


% Set the arabic language as default and english as side langiage
\setmainlanguage[numerals=maghrib]{arabic}
\setotherlanguage{english}

% Set the arabic and the english font
%\newfontfamily{\amiritimes}{Amiri Times}[Script=Arabic]
\newfontfamily\titlefont{DecoType Naskh Swashes}[Script=Arabic]
\newfontfamily\arabicfont[Script=Arabic]{Amiri}
\newfontfamily{\farisifont}[Script=Arabic]{Noto Nastaliq Urdu}
\newfontfamily{\timesfont}{Times New Roman}
\rightfootnoterule
\renewcommand{\footnotesize}{\small}



\usepackage[T1]{fontenc}
\usepackage{times}

% Aliases for fast inserting arabic or english text
\newcommand{\ar}{\textarabic}
\newcommand{\en}{\textenglish}

% Load the MathTime Proffesional package
\usepackage[lite, subscriptcorrection, zswash]{mtpro2}
\straightbraces

\DeclareMathSymbol{0}{\mathalpha}{operators}{`0}
\DeclareMathSymbol{1}{\mathalpha}{operators}{`1}
\DeclareMathSymbol{2}{\mathalpha}{operators}{`2}
\DeclareMathSymbol{3}{\mathalpha}{operators}{`3}
\DeclareMathSymbol{4}{\mathalpha}{operators}{`4}
\DeclareMathSymbol{5}{\mathalpha}{operators}{`5}
\DeclareMathSymbol{6}{\mathalpha}{operators}{`6}
\DeclareMathSymbol{7}{\mathalpha}{operators}{`7}
\DeclareMathSymbol{8}{\mathalpha}{operators}{`8}
\DeclareMathSymbol{9}{\mathalpha}{operators}{`9}

%\makeatletter
%\renewcommand{\@biblabel}[1]{{$\selectlanguage[numerals=maghrib]{arabic} \mathbf{[#1]}$}}
%\renewcommand{\@makefntext}[1]{%
%	\noindent\makebox[2em][r]{\@thefnmark }#1
%}
%\makeatother

\newcommand{\arabicword}[1]{%
	\ifcase#1
	\or الأول
	\or الثاني
	\or الثالث
	\or الرابع
	\or الخامس 
	\or السادس
    \or السابع
	\or الثامن
	\or التاسع
	\or العاشر
	\else \number#1
	\fi
}

\renewcommand{\thechapter}{\arabicword{\value{chapter}}}

% Change the chapter header format
\makeatletter
\def\@makechapterhead#1{%
	\noindent
	\rule{\textwidth}{2pt}\\[10pt]
	{\huge \textbf{الفصل \thechapter}}\\[20pt]
	{\Huge \textbf{#1}}\\[10pt]
	\rule{\textwidth}{2pt}\\[9em]
%	\thispagestyle{empty}   % No page number
%	\newpage
}

\def\@makeschapterhead#1{
	\noindent
	\rule{\textwidth}{2pt}\\[15pt]
	{\huge \textbf{#1}}\\[10pt]
	\rule{\textwidth}{2pt}\\[5em]
}

\@removefromreset{section}{chapter}

%\def\tagform@#1{\maketag@@@{$(#1)$}}
\makeatother


% Chnage the section header format
\titleformat{\section}[block]
{\bfseries\fontsize{20}{17}\selectfont} 
{\thesection}{0.4em}{}


%\titlespacing{\section}{\parindent}{3em}{5pt}

% Change the subsection header format
\titleformat{\subsection}[block]
{\bfseries\fontsize{16}{17}\selectfont} 
{\thesubsection}{1em}{}  



\renewcommand{\thesection}{{\arabic{section}}.}

\renewcommand{\thesubsection}{{\arabic{chapter} - \arabic{section} - \arabic{subsection}}}

\titlecontents{chapter}[0em]{\large\bfseries\vspace{15pt}}{الفصل \thecontentslabel: }{}{\qquad\contentspage}

\titlecontents{chapter*}[0em]{\bfseries\vspace{15pt}}{}{}{\hfill\normalfont{\contentspage}}

\titlecontents{section}[6em]{\normalfont}{}{}{\qquad\contentspage}


\renewcommand{\thetable}{\arabic{chapter} - \arabic{table}}

\renewcommand{\thefigure}{{\textbf{\arabic{figure}}}}


\newtheoremstyle{mystyle}
{\topsep}
{\topsep}
{}
{}
{\bfseries}
{}
{\newline}  
{#1 #2 \textbf{#3}}

\theoremstyle{mystyle}

\newtheorem{theorem}{مبرهنة}[section]
\newtheorem{definition}{تعريف}[section]
\newtheorem{example}{مثال}[section]
\newtheorem{corollary}[theorem]{نتيجة}
%\newtheorem*{note}{ملاحظة}


\renewcommand{\thetheorem}{{\arabic{chapter} - \arabic{section} - \arabic{theorem}}}
\renewcommand{\thedefinition}{{\arabic{chapter} - \arabic{section} - \arabic{definition}}}
\renewcommand{\theexample}{{\arabic{chapter} - \arabic{section} - \arabic{example}}}
\renewcommand{\thecorollary}{{\arabic{chapter} - \arabic{section} - \arabic{corollary}}}

\newenvironment{myproof}{\noindent\textbf{البرهان}\vspace{5pt}\\ \noindent}{\hfill \(\qedsymbol\)}

\newenvironment{solution}{\noindent\textbf{الحل}\vspace{5pt}\\}{}
\newenvironment{note}{\noindent\textbf{ملاحظة}}{}
\newenvironment{notes}{\noindent\textbf{ملاحظات}\vspace{5pt}\\}{}

\pretocmd{\section}{\Needspace{5\baselineskip}}{}{}
\pretocmd{\subsection}{\Needspace{5\baselineskip}}{}{}
\AtBeginEnvironment{definition}{\Needspace{3\baselineskip}}
\AtBeginEnvironment{theorem}{\Needspace{3\baselineskip}}
\AtBeginEnvironment{example}{\Needspace{3\baselineskip}}
\AtBeginEnvironment{solution}{\Needspace{3\baselineskip}}
\AtBeginEnvironment{note}{\Needspace{3\baselineskip}}
\AtBeginEnvironment{myproof}{\Needspace{3\baselineskip}}

\newcommand{\N}{\mathbb{N}}
\newcommand{\Z}{\mathbb{Z}}
\newcommand{\Q}{\mathbb{Q}}
\newcommand{\R}{\mathbb{R}}
\newcommand{\C}{\mathbb{C}}
\newcommand{\LL}{\mathcal{L}}

% Define new math operators
\DeclareMathOperator{\Log}{Log}
\DeclareMathOperator{\lcm}{lcm}
\DeclareMathOperator{\arcsec}{arcsin}
\DeclareMathOperator{\arccsc}{arccsc}
\DeclareMathOperator{\arccot}{arccot}
\DeclareMathOperator{\re}{Re}
\DeclareMathOperator{\im}{Im}

\renewcommand{\mathbf}[1]{\textbf{\en{#1}}}




\begin{document}
	\renewcommand{\thefootnote}{\ar{(\arabic{footnote})}}
	\renewcommand{\figurename}{الشكل}
	\abovedisplayskip=3pt
	\belowdisplayskip=3pt
	\renewcommand{\jot}{10pt}
	\renewcommand{\theequation}{\arabic{equation}.\arabic{chapter}}
%		\arabicfont
	%\begin{titlepage}
	\titlefont
	\noindent
	\begin{minipage}{.5\textwidth}
		\begin{center}
			\large
			\textbf{وِزَارَةُ التَّعْلِيمِ الْعَالِي وَالْبَحْثِ الْعِلْمِيِّ}\\
			\textbf{جَامِعَةُ الْبَصْرَةِ}\\
			\textbf{كُلِيَّةُ التَّرْبِيَةِ لِلْعُلُومِ الصَّرْفَةِ}
		\end{center}
	\end{minipage}
	\hfill
	\begin{minipage}{0.2\textwidth}
		\raggedright\color{black!0}\includegraphics[scale=0.21]{$HOME/Programming/LaTeX/template/Title Page/university.png}   
	\end{minipage}
	\vspace*{\fill}
	\fontsize{60pt}{50pt}\selectfont
	\begin{center}
		\textbf{مُقَدِّمَةٌ في التَّحْليلِ العَقَدِيْ}
	\end{center}
	\vspace*{10pt}
	\begin{center}
		\fontsize{30pt}{30pt}\selectfont
		\textbf{تَأْلِيفٌ} \\
		\textbf{سَجَّادُ أَحْمَدَ جَاسِم}
	\end{center}
	\vspace*{\fill}
\end{titlepage}

	%\chapter*{مقدمة}
	التحليل العقدي هو فرع من فروع الرياضيات الذي يختص بدراسة الدوال التي تُعرّف على الأعداد العقدية. حيث أن الأعداد العقدية هي أعداد تتكون من جزئين: جزء حقيقي وآخر تخيلي. 
	
	يعد هذا المجال أساسيًا لفهم سلوك الدوال في المجال العقدي وكيفية تأثير الأجزاء الحقيقية والتخيلية للأعداد العقدية على تلك الدوال. ومن خلال دراسة هذه الدوال، يمكننا فهم خصائص معقدة تتعلق بالتحليل الرياضي، بالإضافة إلى تطبيقات عملية في العديد من المجالات مثل الفيزياء والهندسة.
	
	تعد التطبيقات العملية للتحليل العقدي واسعة جدًا، فهي تستخدم في العديد من الظواهر الطبيعية مثل تفاعلات الجسيمات في فيزياء الكم، وأيضًا في تحليل الدوائر الكهربائية والنماذج الرياضية في الهندسة. يعتبر التحليل العقدي أداة قوية لفهم العديد من الأنظمة الرياضية المعقدة، كما أنه يعد ركيزة أساسية لحل المعادلات التفاضلية التي تظهر في التطبيقات الفيزيائية والهندسية.
	
	من خلال هذا المجال، يتمكن الباحثون من دراسة التكامل العقدي ودوال أخرى في المستوى العقدي، مما يعزز القدرة على التعامل مع النماذج الرياضية المعقدة وفهم الديناميكيات الدقيقة التي تحكم الأنظمة المختلفة.
	

	%\tableofcontents
	\chapter{الاعداد المرڪبة}

\section{الجمع والضرب} 
\textbf{\textit{الاعداد المرڪبة}} تعرف بأنها الازواج المرتبة $(x, y)$ من الاعداد الحقيقية والتي نتعامل معها ڪنقاط في المستوي.
{
		\begin{tikzpicture}
	\draw[thick, -latex] (0, -0.5) -- (0,4) node[left] {$y$};
	\draw[thick, -latex] (-0.5, 0) -- (6, 0) node[below] {$x$};
	\node[below left] (0, 0) {$O$};
		
		% تعريف الإحداثيات
		\coordinate (Z) at (4,2);
		\coordinate (I) at (0,1);
		\coordinate (X) at (4,0);
		
		% رسم النقاط
		\filldraw[thick, black] (Z) circle (2pt);
		\filldraw[thick, black] (I) circle (2pt);
		\filldraw[thick, black] (X) circle (2pt);
		
		% رسم التسميات بعيداً عن النقاط
		\node[right] at (Z) {$z=(x,y)$};
		\node[left]  at (I) {$i=(0,1)$};
		\node[below] at (X) {$x=(x,0)$};
		
\end{tikzpicture}
\FIGURE{}}


\section*{تمارين}
\begin{enumerate}
	\item تحقق من\\
	\begin{tasks*}
		\item $(1+i)(1-i) = 2$
		\item $(2+i)(2-i) = 5$
		\item $(3+2i)(1-2i) = 7-4i$
	\end{tasks*}
	
	\item اثبت ان\\
	\begin{tasks*}
		\item $\re(i z) = - \im z$
		\item $\im(i z) = \re z$
		\item $\re(z^2) = x^2 - y^2$ من اجل $z = x+iy$
		\item $\im(z^2) = 2xy$ من اجل $z = x+iy$
	\end{tasks*}
	
	\item برهن ان 
	\[
	(1 + z)^3 = 1 + 3z + 3z^2 + z^3
	\]
	
	\item اثبت ان كلا العددين 
	\[
	z = i, \quad z = -i
	\] 
	يحققان المعادلة $z^2 + 1 = 0$.
	
	\item حل المعادلة 
	\[
	z^2 - 2z + 5 = 0
	\] 
	من اجل $z = (x, y)$ من خلال كتابة
	\[
	(x, y)(x, y) - 2(x, y) + (5, 0) = (0, 0)
	\]
	ومن ثم حل نظام من المعادلات في $x$ و $y$.
	
	\SUG{استخدم التمثيل البياني للعدد المركب أو حل المعادلة كنظام: $x^2 - y^2 - 2x + 5 = 0$, $2xy - 2y = 0$}
	
	\ANS{$z = (1, \pm 2)$}
	
	\item حل المعادلة 
	\[
	z^3 = 1
	\] 
	واحصل على جميع الجذور من الشكل $z = (x, y)$.
	
	\ANS{
		\[
		z_1 = (1,0), \quad z_2 = \left(-\dfrac{1}{2}, \dfrac{\sqrt{3}}{2}\right), \quad z_3 = \left(-\dfrac{1}{2}, -\dfrac{\sqrt{3}}{2}\right)
		\]
	}
	
	\item \begin{tasks*}
		\item اوجد الناتج: $(2+3i) + (1-4i)$
		\ANS{$3 - i$}
		
		\item اوجد الناتج: $(1+i)(1-i)$
		\ANS{$2$}
		
		\item اوجد الناتج: $(3+i)^2$
		\ANS{$8 + 6i$}
		
		\item اوجد الناتج: $\dfrac{2+3i}{1-i}$
		\ANS{$\dfrac{-1+5i}{2}$}
	\end{tasks*}
\end{enumerate}

\section{بعض الخواص الجبرية}

\section*{تــماريــن}
\begin{enumerate}
	\item بسط المقادير التالية الى ابسط صورة\\
	\begin{tasks}
		\item $\dfrac{1+2i}{3-4i} + \dfrac{2-i}{5i}$ ؛ 
		\item $\dfrac{5i}{(1-i)(2-i)(3-i)}$ ؛ 
		\item $(1-i)^4$
	\end{tasks}\\[10pt]
	\ANS{
	\begin{tasks}
		\item $-2/5$ ؛
		\item $-1/2$ ؛ 
		\item $-4$
	\end{tasks}
	}
	\item اثبت أن
	\[
	\frac{1}{1/z} = z \qquad (z\neq 0)
	\]
	
	\item استعمل قواعد التجميع والتوزيع لعملية الضرب لتبرهن أن
	\[
	(z_1 z_2)(z_3 z_4) = (z_1 z_3)(z_2 z_4)
	\]
\end{enumerate}

\section{المزيد من الخواص الجبرية}

\section{المتجهات والمعيار}



{\noindent
\begin{tikzpicture}[scale=2]
		\draw[thick] (-2.5, 0) -- (2.5, 0) node[below] {$x$};
		\draw[thick] (0, -0.25) -- (0, 1.5) node[left] {$y$};

		\node[below left] at (0, 0) {$O$};
		
		\draw[thick] (0, 1) -- (-0.1, 1) node[left] {$1$};
		\draw[thick] (-2, 0) -- (-2, -0.1) node[below] {$-2$};
		
		\draw[very thick, -latex, shorten >=4pt] (0, 0) -- node[sloped, above] {$z=x+iy$} (1.5, 0.8);
		
		\draw[fill=black] (1.5, 0.8) circle (1pt) node[right] {$z=(x, y)$};
		
		\draw[very thick, -latex, shorten >=4pt] (0, 0) -- node[sloped, above] {$-2+i$} (-2, 1);
		
		\draw[fill=black] (-2, 1) circle (1pt) node[above] {$(-2,1)$};
\end{tikzpicture}
\FIGURE{}}

{\noindent
\begin{tikzpicture}[scale=2.5]
		\draw[thick] (-0.25, 0) -- (2.6, 0) node[below] {$x$};
		\draw[thick] (0, -0.25) -- (0, 1.8) node[left] {$y$};
		
		\node[below left] at (0, 0) {$O$};
		
		\draw[very thick, -latex, shorten >= 4pt] (0, 0) -- node[left, pos=0.65] {$z_2$} (0.5, 1);
		\draw[very thick, -latex, shorten >= 4pt] (0, 0) -- node[below, pos=0.65]{$z_1$} (1.5, 0.5);
		\draw[very thick, -latex, shorten >= 4pt] (0, 0) --node[sloped, above]{$z_1+z_2$} (2, 1.5);
		
		\draw[fill=black] (0.5, 1) circle (1pt);
		\draw[fill=black] (1.5, 0.5) circle (1pt);
		\draw[fill=black] (2, 1.5) circle (1pt);
		
		\draw[mydash, thick, shorten >=14pt] (0.5, 1) -- (2, 1.5);
		\draw[very thick, mydash, -latex, shorten >=6pt] (1.5, 0.5) -- (2, 1.5);
\end{tikzpicture}
\FIGURE{}}


{\noindent
\begin{tikzpicture}[scale=2.5]
		\draw[thick] (-0.25, 0) -- (2.6, 0) node[below] {$x$};
		\draw[thick] (0, -0.25) -- (0, 1.8) node[left] {$y$};
		
		\node[below left] at (0, 0) {$O$};
		
		\coordinate (O) at (0, 0);
		\coordinate (z1) at (2, 0.5);
		\coordinate (z2) at (0.5, 1);
		\coordinate (z3) at (1.5, -0.5);
		\draw[very thick, -latex, shorten >= 4pt] (O) -- node[left, pos=0.65] {$z_2$} (z2);
		\draw[very thick, -latex, shorten >= 4pt] (O) -- node[below, pos=0.65]{$z_1$} (z1);
		
		
		\draw[fill=black] (z2) circle (1pt) node[above]{$(x_2, y_2)$};
		\draw[fill=black] (z1) circle (1pt) node[right]{$(x_1, y_1)$};
		\draw[fill=black] (z3) circle (1pt);
		
		\draw[mydash, very thick, -latex, shorten >=4pt] (z1) --node[right, pos=0.65]{$-z_2$} (z3);
		\draw[mydash, thick] (z2) --node[sloped, above]{$|z_1-z_2|$} (z1);
		
		\draw[very thick, -latex, shorten >=4pt] (O) --node[sloped, below]{$z_1-z_2$} (z3);
		
\end{tikzpicture}
\FIGURE{}}

\section{المتراجحة المثلثية}

\section{المــُرافق العقدي}

{\noindent
\begin{tikzpicture}[scale=2]
		\draw[thick] (-0.5, 0) -- (2.6, 0) node[below] {$x$};
		\draw[thick] (0, -0.5) -- (0, 1.8) node[left] {$y$};
		\node[below left] at (0, 0) {$O$};
		
		\draw[very thick, -latex, shorten >=4pt] (0, 0) --node[above]{$z$} (1.5, 0.5);
		\draw[very thick, -latex, shorten >=4pt] (0, 0) --node[below]{$\conj{z}$} (1.5, -0.5);
		\draw[thick, mydash] (1.5, 0.5) -- (1.5, -0.5);
		
		\draw[fill=black] (1.5, 0.5) circle (1pt) node[right]{$(x, y)$};
		\draw[fill=black] (1.5, -0.5) circle (1pt) node[right]{$(x, -y)$};
\end{tikzpicture}
\FIGURE{}}

\section{الصيغة الأُسية}

{\noindent
\begin{tikzpicture}[scale=1, thick]
	\draw (-pi+1, 0) -- (pi+1, 0) node[above, pos=0.99]{$x$};
	\draw (0, -2) -- (0, 3) node[left, pos=0.96]{$y$};
	\coordinate (A) at (70:2);
	\draw[fill=black] (A) circle (2pt) node[right]{$z = x+iy$};
	\draw[-latex, 
	shorten >=3pt, very thick
	] (0:0) -- (A);
	\draw (0:.3) arc(0:90:.3 and .4);
	\draw (90:.4) arc(90:180:.5 and .4);
	\draw (180:.5) arc(180:270:.5 and .6);
	\draw (270:.6) arc(270:360:.7 and .6);
	\draw[-latex] (0:.7) arc(0:69:.7) node[right, pos=0.55]{$\theta$};
\end{tikzpicture}
\FIGURE{}}

{\noindent
\begin{tikzpicture}[scale=2]
	\draw[thick] (-1.5, 0) -- (1.5, 0) node[below, pos=0.98]{$x$};
	\draw[thick] (0, -1.5) -- (0, 1.5) node[left, pos=0.97]{$y$};
	\node[below left] at (0, 0) {$O$};
	
	\coordinate (A) at (135:1);
	
	\draw[thick] (0:0) circle (1);
	
	\draw[fill=black] (A) circle (1pt) node[left]{$e^{i\theta}$};
	
	\draw[very thick, -latex, shorten >=4pt] (0:0) --node[above]{$1$} (A);
	
	\draw[-latex, thick] (0:0.4) arc(0:135:0.4) node[pos=0.23, right]{$\theta$};
\end{tikzpicture}
\FIGURE{}}

\section{الضرب والقوى في الصيغة الأسية}

\section{السعة للضرب والقسمة}

{\noindent
\begin{tikzpicture}[scale=2,thick]
	\draw[thick] (-.5,0)--(3.5,0) node[below]{$x$};
	\draw[thick] (0,-.5)--(0,3) node[left]{$y$};
	\def\aone{15} \def\atwo{45}
	\def\rone{2}  \def\rtwo{1.6}
	\draw[fill=black] (\aone:\rone) circle (1pt) node[right] {$z_1^{}$};
	\draw[fill=black] (\atwo:\rtwo) circle (1pt) node[right] {$z_2^{}$};
	\draw[fill=black] (\aone+\atwo:\rone*\rtwo) circle (1pt) node[right] {$z_1^{}z_2^{}$};
	\draw[very thick, -latex, shorten >=4pt]
	(0,0)--(\aone:\rone);
	\draw[very thick, -latex, shorten >=4pt] (0,0)--(\atwo:\rtwo);
	\draw[very thick, -latex, shorten >=4pt] (0,0)--(\aone+\atwo:\rone*\rtwo);
	\draw[-latex] (0:1.5) arc(0:\aone:1.5) 
	node[pos=.5,right]{$\theta_1^{}$};
	\draw[-latex] (0:.8) arc(0:\atwo:.8) 
	node[pos=.75,right]{$\theta_2^{}$};
	\draw[-latex] (0:2.5) arc(0:\aone+\atwo:2.5)
	node[pos=.5,right]{$\theta_1^{}+\theta_2^{}$};
\end{tikzpicture}
\FIGURE{}}

\section{جذور الاعداد العقدية}

{\noindent
\begin{tikzpicture}[thick]
	\draw[] (-2.5, 0) -- (4, 0) node[below]{$x$};
	\draw[] (0, -2.5) -- (0, 3) node[left]{$y$};
	\draw[] (0, 0) circle (2);
	\node[below left] at (0, 0) {$O$};
	\node[below right] at (0:2) {$\sqrt[n]{r_0}$};
	\coordinate (A) at (110:2);
	\coordinate (B) at (160:2);
	\draw[fill=black] (A) circle (2pt) node[above left]{$c_{k-1}$};
	\draw[fill=black] (B) circle (2pt) node[left]{$c_k$};
	\draw[very thick, -latex, shorten >=3pt] (0, 0) -- (A);
	\draw[very thick, -latex, shorten >=3pt] (0, 0) -- (B);
	\draw[-latex] (110:0.7) arc(110:155:0.7) node[pos=0.5, above left]{$\frac{2\pi}{n}$};
\end{tikzpicture}
\FIGURE{}}

{\noindent
\begin{tikzpicture}[scale=2, thick]
	\coordinate (A) at (45:1);
	\coordinate (B) at (135:1);
	\coordinate (C) at (225:1);
	\coordinate (D) at (315:1);
	\draw[] (-1.5, 0) -- (1.5, 0) node[below]{$x$};
	\draw[] (0, -1.5) -- (0, 1.5) node[left]{$y$};
	\node[below right] at (0:1) {$z$};
	\draw[] (0, 0) circle (1);
	\draw[fill=black] (A) circle (1pt) node[right]{$c_0^{}$};
	\draw[fill=black] (B) circle (1pt) node[left]{$c_1^{}$};
	\draw[fill=black] (C) circle (1pt) node[left]{$c_2^{}$};
	\draw[fill=black] (D) circle (1pt) node[right]{$c_3^{}$};
	
	\draw[very thick, -latex, shorten >=3pt] (0, 0) -- (A);
	\draw[very thick, -latex, shorten >=3pt] (0, 0) -- (B);
	\draw[very thick, -latex, shorten >=3pt] (0, 0) -- (C);
	\draw[very thick, -latex, shorten >=3pt] (0, 0) -- (D);
	
	\draw[mydash] (A) -- (B);
	\draw[mydash] (B) -- (C);
	\draw[mydash] (C) -- (D);
	\draw[mydash] (D) -- (A);
\end{tikzpicture}
\FIGURE{}}

{\noindent
	\begin{tikzpicture}[scale=2, thick, baseline=0]
	\draw (-1.3, 0) -- (1.3, 0) node[below]{$x$};
	\draw (0, -1.3) -- (0, 1.3) node[left]{$y$};
	
	\coordinate (A) at (0:1);
	\coordinate (B) at (120:1);
	\coordinate (C) at (240:1);
	\draw[fill=black] (A) circle (1pt)  node[below right]{$1$};
	\draw[fill=black] (B) circle (1pt)  node[above left]{$\omega_3^{}$};
	\draw[fill=black] (C) circle (1pt) node[below left]{$\omega_3^2$};
	
	\draw (0, 0) circle (1);
	\draw[-latex, very thick, shorten >=3pt] (0, 0) -- (A);
	\draw[-latex, very thick, shorten >=3pt] (0, 0) -- (B);
	\draw[-latex, very thick, shorten >=3pt] (0, 0) -- (C);
	
	\draw[mydash] (A) -- (B) -- (C) -- (A);
\end{tikzpicture}
\begin{tikzpicture}[scale=2, thick, baseline=0]
	\draw (-1.3, 0) -- (1.3, 0) node[below]{$x$};
	\draw (0, -1.3) -- (0, 1.3) node[left]{$y$};
	
	\coordinate (A) at (0:1);
	\coordinate (B) at (90:1);
	\coordinate (C) at (180:1);
	\coordinate (D) at (270:1);
	\draw[fill=black] (A) circle (1pt)  node[below right]{$1$};
	\draw[fill=black] (B) circle (1pt)  node[above right]{$\omega_4^{}$};
	\draw[fill=black] (C) circle (1pt) node[above left]{$\omega_4^2$};
	\draw[fill=black] (D) circle (1pt) node[below left]{$\omega_4^3$};
	\draw (0, 0) circle (1);
	\draw[-latex, very thick, shorten >=3pt] (0, 0) -- (A);
	\draw[-latex, very thick, shorten >=3pt] (0, 0) -- (B);
	\draw[-latex, very thick, shorten >=3pt] (0, 0) -- (C);
	\draw[-latex, very thick, shorten >=3pt] (0, 0) -- (D);
	\draw[mydash] (A) -- (B) -- (C) -- (D) -- (A);
\end{tikzpicture}
\begin{tikzpicture}[scale=2, thick, baseline=0]
	\draw (-1.3, 0) -- (1.3, 0) node[below]{$x$};
	\draw (0, -1.3) -- (0, 1.3) node[left]{$y$};
	
	\coordinate (A) at (0:1);
	\coordinate (B) at (60:1);
	\coordinate (C) at (120:1);
	\coordinate (D) at (180:1);
	\coordinate (E) at (240:1);
	\coordinate (F) at (300:1);
	\draw[fill=black] (A) circle (1pt)  node[below right]{$1$};
	\draw[fill=black] (B) circle (1pt)  node[above right]{$\omega_6^{}$};
	\draw[fill=black] (C) circle (1pt) node[above left]{$\omega_6^2$};
	\draw[fill=black] (D) circle (1pt) node[above left]{$\omega_6^3$};
	\draw[fill=black] (E) circle (1pt) node[below left]{$\omega_6^4$};
	\draw[fill=black] (F) circle (1pt) node[below right]{$\omega_6^5$};
	\draw (0, 0) circle (1);
	\draw[-latex, very thick, shorten >=3pt] (0, 0) -- (A);
	\draw[-latex, very thick, shorten >=3pt] (0, 0) -- (B);
	\draw[-latex, very thick, shorten >=3pt] (0, 0) -- (C);
	\draw[-latex, very thick, shorten >=3pt] (0, 0) -- (D);
	\draw[-latex, very thick, shorten >=3pt] (0, 0) -- (E);
	\draw[-latex, very thick, shorten >=3pt] (0, 0) -- (F);
	\draw[mydash] (A) -- (B) -- (C) -- (D) -- (E) -- (F) -- (A);
\end{tikzpicture}\\
\FIGURE{}}

\section{المناطق في المستوي العقدي}

{\noindent
\begin{tikzpicture}[thick]
	\draw[mydash, fill=gray!50] (0, 0) circle (2);
	\draw[mydash, fill=white] (0, 0) circle (1);
	
	\draw (-3, 0) -- (3.5, 0) node[below, pos=0.99]{$x$};
	\draw (0, -3) -- (0, 3) node[left, pos=0.96]{$y$};
	\node[below left] at (0, 0) {$O$};
	\node[below right] at (1, 0) {$1$};
	\node[below right] at (2, 0) {$2$};
	
	\coordinate (z1) at (45:1.4);
	\coordinate (z2) at (200:1.7);
	
	\draw[fill=black] (z1) circle (2pt) node[right]{$z_2$};
	\draw[fill=black] (z2) circle (2pt) node[right]{$z_1$};
	\draw (z1) -- (110:1.7) -- (z2);
\end{tikzpicture}
\FIGURE{}}

{\noindent
\begin{tikzpicture}[scale=3, thick]
	\draw (-1, 0) -- (1, 0) node[below, pos=0.99]{$x$};
	\draw (0, -1.3) -- (0, 1/3) node[left, pos=0.96]{$y$};
	\node[below left] at (0, 0) {$O$};
	\draw[mydash] (0, -1/2) circle (1/2);
	\draw[fill=black] (0, -1/2) circle (2/3pt) node[right]{$-\dfrac{i}{2}$};
\end{tikzpicture}
\FIGURE{}}









	\chapter{الدوال التحليلية}

\section{دوال وتحويلات}

\section*{تمارين على التكاملات}
\begin{enumerate}
	\item تحقق من \\ 
	\begin{tasks*}
		\item $\displaystyle \int 1\,dx = x + C$ 
		\ANS{$x + C$}
		\item $\displaystyle \int x\,dx = \frac{x^2}{2} + C$
		\ANS{$\frac{x^2}{2} + C$}
		\item $\displaystyle \int x^2\,dx = \frac{x^3}{3} + C$
		\ANS{$\frac{x^3}{3} + C$}
	\end{tasks*}
	
	\item احسب التكاملات التالية \\ 
	\begin{tasks*}
		\item $\displaystyle \int (2x + 3)\,dx$
		\ANS{$x^2 + 3x + C$}
		\item $\displaystyle \int (x^2 - 5x + 4)\,dx$
		\ANS{$\frac{x^3}{3} - \frac{5x^2}{2} + 4x + C$}
		\item $\displaystyle \int (3\cos x + 4\sin x)\,dx$
		\ANS{$3\sin x - 4\cos x + C$}
	\end{tasks*}
	
	\item برهن ان \\ 
	\begin{tasks*}
		\item $\displaystyle \int \frac{d}{dx}(x^n) dx = x^n + C$
		\item $\displaystyle \int \frac{d}{dx}(\sin x) dx = \sin x + C$
		\item $\displaystyle \int \frac{d}{dx}(\cos x) dx = \cos x + C$
	\end{tasks*}
	
	\item استخدم التعويض \\ 
	\begin{tasks*}
		\item $\displaystyle \int (2x+1)^5 dx$
		\SUG{ضع $u = 2x+1$ ثم $du=2dx$}
		\ANS{$\frac{(2x+1)^6}{12} + C$}
		
		\item $\displaystyle \int \sin(3x) dx$
		\SUG{ضع $u = 3x$ ثم $du = 3 dx$}
		\ANS{$-\frac{1}{3} \cos(3x) + C$}
		
		\item $\displaystyle \int x e^{x^2} dx$
		\SUG{ضع $u = x^2$ ثم $du = 2x dx$}
		\ANS{$\frac{1}{2} e^{x^2} + C$}
	\end{tasks*}
	
	\item التكاملات الجزئية \\ 
	\begin{tasks*}
		\item $\displaystyle \int x \cos x\, dx$\\
		\SUG{استخدم التكامل بالجزء: $u=x$, $dv=\cos x dx$}\\
		\ANS{$x \sin x + \cos x + C$}
		
		\item $\displaystyle \int x e^x dx$
		\SUG{التكامل بالجزء: $u=x$, $dv=e^x dx$}
		\ANS{$(x-1)e^x + C$}
	\end{tasks*}
	
	\item تكاملات المشتقات \\ 
	\begin{tasks*}
		\item $\displaystyle \int \frac{1}{x} dx$
		\ANS{$\ln |x| + C$}
		\item $\displaystyle \int \frac{1}{1+x^2} dx$
		\ANS{$\arctan x + C$}
		\item $\displaystyle \int \frac{1}{\sqrt{1-x^2}} dx$
		\ANS{$\arcsin x + C$}
	\end{tasks*}
	
	\item تكاملات متعددة الحدود \\ 
	\begin{tasks*}
		\item $\displaystyle \int (3x^3 - 2x^2 + x - 5) dx$
		\ANS{$\frac{3x^4}{4} - \frac{2x^3}{3} + \frac{x^2}{2} - 5x + C$}
		\item $\displaystyle \int (x^4 + 2x^2 + 1) dx$
		\ANS{$\frac{x^5}{5} + \frac{2x^3}{3} + x + C$}
		\item $\displaystyle \int (x^3 - 3x + 2) dx$
		\ANS{$\frac{x^4}{4} - \frac{3x^2}{2} + 2x + C$}
	\end{tasks*}
\end{enumerate}

\section[التحويل $w = z^2$]{التحويل \en{\textit{w = z}\textsuperscript{2}}}

{\noindent
\begin{tikzpicture}[thick,  baseline=-3.5cm]
	\draw (-3.5, 0) -- (3.5, 0) node[below]{$x$};
	\draw (0, -3.5) -- (0, 3.5) node[left]{$y$};
	\node[below left] at (0, 0) {$O$};
	\draw[mydash, very thick, arrow inside={0.45}{for}] plot[domain=1/6:3, samples=100]({\x}, {1/(2*\x)});
	\draw[very thick, mydash, arrow inside={0.45}{for}] plot[domain=-1/6:-3, samples=100]({\x}, {1/(2*\x)});
	
	\draw[very thick, arrow inside={0.3}{for}] plot[domain=1.5:-1.5, samples=100]({-cosh(\x)}, {sinh(\x)}); 
	\draw[very thick, arrow inside={0.3}{for}] plot[domain=1.5:-1.5, samples=100]({cosh(\x)}, {-sinh(\x)});
\end{tikzpicture}\qquad
\begin{tikzpicture}[thick,  baseline=-3.5cm]
	\draw (-3.5, 0) -- (3.5, 0) node[below, pos=0.99]{$u$};
	\draw (0, -3.5) -- (0, 3.5) node[left, pos=0.96]{$v$};
	\node[below left] at (0, 0) {$O$};
	
	\draw[mydash, very thick, arrow inside={0.4}{for}] (-3.5, 1.5) -- (2, 1.5) node[right]{$v = c_2 > 0$};
	
	\draw[very thick, arrow inside={0.4}{for}] (0.5, -3.5) -- (.5, 2.5) node[above right]{$u=c_1>0$};
\end{tikzpicture}
\FIGURE{$w = z^2$}}

{\noindent
\begin{tikzpicture}[thick, baseline=-0.5cm]
	\draw (-0.5, 0) -- (3, 0) node[below]{$x$};
	\draw (0, -.5) -- (0, 3) node[left]{$y$};
	\node[below left] at (0, 0) {$O$};
	\draw[very thick] (0, 0) -- (0, 3)
	(0, 0) -- (3, 0)
	;
	\coordinate (A) at (0:2);
	\draw[fill=black] (A) circle (2pt) node[below]{$r_0^{}$};
	\draw[mydash, arrow inside={0.9}{for}] (0, 0) -- (30:3.3);
	\draw[arrow inside={0.5}{for}] (A) arc(0:90:2);
\end{tikzpicture}\qquad\qquad
\begin{tikzpicture}[thick, baseline=-0.5cm]
	\draw[very thick] (-3.5, 0) -- (3.5, 0) node[below]{$u$};
	\draw (0, -.5) -- (0, 3.5) node[left]{$v$};
	\node[below left] at (0, 0) {$O$};
	\draw[mydash, arrow inside={0.9}{for}] (0, 0) -- (60:3.3);
	\coordinate (A) at (0:2.5);
	\draw[fill=black] (A) circle (2pt) node[below]{$r_0^{2}$};
	\draw[arrow inside={0.75}{for}] (A) arc(0:180:2.5);
\end{tikzpicture}
\FIGURE{$w = z^2$}}

\section{الغايات}

{
\noindent
\begin{tikzpicture}[thick, baseline=-.5cm]
	\draw (-0.5, 0) -- (3, 0) node[below]{$x$};
	\draw (0, -.5) -- (0, 3) node[left]{$y$};
	\node[below left] at (0, 0) {$O$};
	
	\coordinate (z0) at (1.5, 0.5);
	\coordinate (z) at (1.9, 0.25);
	
	\draw[mydash] (z0) circle (1);
	\draw[fill=black] (z) circle (2pt) node[right]{$z$}; 
	\draw (z0) --node[above, pos=0.25]{$\delta$} +(45:1);
	\draw[fill=white] (z0) circle (2pt) node[left]{$z_0^{}$}; 
\end{tikzpicture}\qquad\qquad
\begin{tikzpicture}[thick, baseline=-.5cm]
	\draw (-3, 0) -- (0.5, 0) node[below]{$u$};
	\draw (0, -.5) -- (0, 3) node[left]{$v$};
	\node[below left] at (0, 0) {$O$};
	
	\coordinate (z0) at (-1.5, 1.5);
	\coordinate (z) at (-1, 1);
	
	\draw[mydash] (z0) circle (1.25);
	\draw[fill=black] (z) circle (2pt) node[left]{$w$}; 
	\draw (z0) --node[above, pos=0.25]{$\varepsilon$} +(45:1.25);
	\draw[fill=white] (z0) circle (2pt) node[left]{$w_0^{}$}; 
\end{tikzpicture}
\FIGURE{}
}

	\chapter{التڪامل العقدي}

{\noindent
	\begin{tikzpicture}[scale=2.5, thick]
		\draw (-0.3, 0) -- (2.5, 0) node[above, pos=0.99]{$x$};
		\draw (0, -0.3) -- (0, 1.5) node[left, pos=0.96]{$y$};
		\coordinate (A) at (0, 0);
		\coordinate (B) at (1, 1);
		\coordinate (C) at (2, 1);
		\draw[fill=black] (A) circle (1.25pt) node[below left]{$O$};
		\draw[fill=black] (B) circle (1.25pt) node[above]{$1+i$};
		\draw[fill=black] (C) circle (1.25pt) node[above]{$2+i$};
		\draw[thick] (0, 1) -- (-0.07, 1) node[left]{$1$};
		\draw (1, 0) -- (1, -.07) node[below]{$1$};
		\draw (2, 0) -- (2, -.07) node[below]{$2$};
		\draw[very thick, arrow inside={0.6}{for}] (A) -- (B);
		\draw[very thick, arrow inside={0.6}{for}] (B) -- (C);
	\end{tikzpicture}
\FIGURE{}}

{\noindent
\begin{tikzpicture}[thick, scale=2.5]
	\draw (-0.3, 0) -- (2.5, 0) node[above, pos=0.99]{$x$};
	\draw (0, -0.3) -- (0, 1.5) node[left, pos=0.96]{$y$};
	\node[below left] at (0, 0) {$O$};
	
	\coordinate (z1) at (0.4, 0.4);
	\coordinate (z2) at (2, 1);
	
	\draw[fill=black] (z1) circle (1pt) node[below]{$z_1^{}$};
	\draw[fill=black] (z2) circle (1pt) node[below]{$z_2^{}$};
	\draw[very thick,
	arrow inside={0.25}{rev},
	arrow inside={0.7}{for}] (z1) to[bend left=20]node[pos=0.7, above]{$C$} node[pos=.25, above, xshift=-8pt]{$-C$} (z2);

\end{tikzpicture}
\FIGURE{}}

{\noindent
\begin{tikzpicture}[thick, scale=2]
	\draw (-0.3, 0) -- (4, 0) node[above, pos=0.99]{$x$};
	\draw (0, -0.3) -- (0, 2) node[left, pos=0.96]{$y$};
	\node[below left] at (0, 0) {$O$};
	
	\coordinate (z1) at (0.5, 0.5);
	\coordinate	(z2) at (1, 1.5);
	\coordinate(z3) at (3, 1);
	
	\draw[fill=black] (z1) circle (1.25pt)node[below]{$z_1^{}$};
	\draw[fill=black] (z2) circle (1.25pt)node[above]{$z_2^{}$};
	\draw[fill=black] (z3) circle (1.25pt)node[below]{$z_3^{}$};
	
	\draw[very thick, arrow inside={0.5}{for}] (z1) to [bend left]node[pos=0.5, above left]{$C_1$} (z2);
	\draw[very thick, arrow inside={0.5}{for}] (z2) to [bend right]node[pos=0.5, above]{$C_2$} (z3);
	
	\node at (1.25, 0.7) {$C$};
\end{tikzpicture}
\FIGURE{}}
	
	\textbf{مثال 2.} العدد $-1 - i$ في المثال 1 له الصيغة الأسية التالية:
	\begin{equation}
		-1 - i = \sqrt{2} \exp \left[ i \left( -\frac{3\pi}{4} \right) \right]. \tag{5}
	\end{equation}
	
	مع الاتفاق على أن $e^{-i\theta} = e^{i(-\theta)}$، يمكن أيضاً كتابة هذا التعبير كـ $-1 - i = \sqrt{2} e^{-i3\pi/4}$. التعبير (5) هو بالطبع مجرد واحد من عدد لا نهائي من الاحتمالات للصيغة الأسية للعدد $-1 - i$:
	\begin{equation}
		-1 - i = \sqrt{2} \exp \left[ i \left( -\frac{3\pi}{4} + 2n\pi \right) \right] \quad (n = 0, \pm 1, \pm 2, \dots). \tag{6}
	\end{equation}
	
	لاحظ كيف يخبرنا التعبير (4) عندما تكون $r = 1$ أن الأعداد $e^{i\theta}$ تقع على الدائرة التي مركزها نقطة الأصل ونصف قطرها الوحدة، كما هو موضح في الشكل 7. قيم $e^{i\theta}$ تكون حينها مباشرة من ذلك الشكل، دون الرجوع إلى صيغة أويلر. فمن الواضح هندسياً، على سبيل المثال، أن:
	\[
	e^{i\pi} = -1, \quad e^{-i\pi/2} = -i, \quad \text{\textit{and}} \quad e^{-i4\pi} = 1.
	\]
	
	\begin{thebibliography}{9}
	\bibitem{key1} الدڪتور احمد الاسدي، مقدمة في التحليل العقدي، جامعة البصرة، ١٩٩٦.
	\bibitem{key2} الدكتورة سارة الجبوري، مدخل إلى نظرية الأعداد، جامعة بغداد، ٢٠٠١.
	\bibitem{key3} الدكتور مصطفى الخطيب، تحليل البيانات باستخدام الخوارزميات، جامعة الكوفة، ٢٠٠٨.
	\bibitem{key4} الدكتور خالد الشمري، الميكانيكا الكلاسيكية وتطبيقاتها، جامعة الموصل، ٢٠١٢.
	\bibitem{key5} الدكتور عماد يوسف، مقدمة في الجبر الخطي، جامعة الفاتح، ٢٠١٥.
	\bibitem{key6} الدكتور يوسف القيسي، نظرية المجموعات والتطبيقات، جامعة تكريت، ٢٠١٨.
	
	\setLTR\timesfont
		\bibitem{key1} Dr. John Smith, Introduction to Complex Analysis, University of Oxford, 1996.
		\bibitem{key2} Dr. Emily Johnson, Fundamentals of Number Theory, Harvard University, 2001.
		\bibitem{key3} Dr. Michael Brown, Data Analysis with Algorithms, Stanford University, 2008.
		\bibitem{key4} Dr. David Wilson, Classical Mechanics and Applications, University of Cambridge, 2012.
		\bibitem{key5} Dr. Sarah Taylor, Introduction to Linear Algebra, MIT, 2015.
		\bibitem{key6} Dr. William Clark, Set Theory and Applications, University of California, Berkeley, 2018.
	
\end{thebibliography}

\end{document}